% Programming with InitechBase, Chapter 6: Communicating with the User
% (AN-16; annex Sec 8.2 P6).  Panels replayed (examples/p6-*.dot); sources
% cmd/proc.c (ACCEPT, INPUT, WAIT), cmd/query.c (? and ??).
\chapter{Communicating with the User}
\Gates{test-pubs-examples, test-interp-proc}
A program that does the same thing every time needs no questions.  Most need to
know something first: which invoice, which customer, which choice of the menu.
This chapter shows how a program asks, and how it answers.

\begin{ChapterCovers}
\item Displaying results: \cmd{?} and \cmd{??}
\item Asking for characters: \cmd{ACCEPT}
\item Asking for a value: \cmd{INPUT}
\item Asking for a key: \cmd{WAIT}
\item Finding an invoice by its number
\end{ChapterCovers}

\ATEHead{Preparing for This Chapter}
Start InitechBase with a fresh copy of the sample diskette, and type
\type{USE INVOICE INDEX INVNO}.

\ATEHead{Displaying Results}
\cmd{?} begins a new line and displays its values; \cmd{??} displays them where
the cursor is.  A program builds a line from several \cmd{??} commands, or from
one \cmd{?} and the \texttt{+} operator.  \cmd{?} with nothing after it prints an
empty line, which a program uses to separate its output.

\ATEHead{Asking for Characters}
\cmd{ACCEPT} displays a prompt and stores the line the user types in a memory
variable, as characters:

\SyntaxLine{\kw{ACCEPT} \xp{prompt} \kw{TO} \xp{memvar}}

The user need not type quotation marks.  Convert the answer to capitals with
\fn{UPPER()} before you compare it with fields kept in capitals.

\ATEHead{Asking for a Value}
\cmd{INPUT} evaluates the line the user types and stores its value, of whatever
type: a number for an amount, \texttt{.T.} or \texttt{.F.} for a yes or no.  A
character string must then be typed in quotation marks, so the system uses
\cmd{ACCEPT} for characters and \cmd{INPUT} for amounts.

\ATEHead{Asking for a Key}
\cmd{WAIT} takes one key, without \key{Enter}.  The menu of Chapter~11 uses it
for the clerk's choice.  \cmd{WAIT} without \kw{TO} only waits, to hold output
on the screen.

\ATEHead{Finding an Invoice}
\file{FINDINV.PRG} asks for an invoice number and displays the invoice:
\ProgramListing{FINDINV.PRG}
\begin{Steps}
\item Type \type{DO FINDINV}.
\ATEResponse{The program asks \texttt{Invoice number:}.}
\item Type \type{a1002} and press \key{Enter}.
\item Type \type{DO FINDINV} again and answer \type{A9999}.
\end{Steps}
\DotExample{examples/p6-findinv}
The answer is converted to capitals and looked up in the index \file{INVNO.NDX}
with \cmd{SEEK}.  \fn{FOUND()} decides which of the two messages the program
prints.

\ATEHead{Checking an Answer}
A program cannot trust what is typed.  \file{AMTCHECK.PRG} asks for an amount
until the answer is a number in the range it accepts:
\ProgramListing{AMTCHECK.PRG}
\DotExample{examples/p6-amtcheck}
\fn{TYPE()} tells the program whether the answer is a number before it is
compared; a character string typed in quotation marks is refused like a number
out of range.  The loop ends only when \texttt{MOK} is true.

\ATEHead{Where Do You Go From Here?}
Chapter~7 adds records to the ledger from the answers to questions.

\ATEHead{Summary}
\begin{enumerate}
\item \cmd{?} and \cmd{??} display; \cmd{?} alone prints an empty line.
\item \cmd{ACCEPT} asks for characters; \cmd{INPUT} for a value of any type;
  \cmd{WAIT} for one key.
\item Convert character answers with \fn{UPPER()} before comparing them.
\end{enumerate}
